\begin{problem}[第31届全国中学生物理竞赛复赛][圆盘成像与光阑]{21}
如图所示，现有一圆盘状发光体，其半径为 $5\mathrm{cm}$，放置在一焦距为 $10\mathrm{cm}$、半径为 $15\mathrm{cm}$ 的凸透镜前，圆盘与凸透镜的距离为 $20\mathrm{cm}$，透镜后放置一半径大小可调的圆形光阑和一个接收圆盘像的光屏。图中所有光学元件相对于光轴对称放置。请在几何光学近轴范围内考虑下列问题，并忽略像差和衍射效应。
    \begin{subproblem}
        未放置圆形光阑时，给出圆盘像的位置、大小、形状；
    \end{subproblem}
    \begin{subproblem}
        若将圆形光阑放置于凸透镜后方 $6\mathrm{cm}$ 处。当圆形光阑的半径逐渐减小时，圆盘的像会有什么变化？是否存在某一光阑半径 $r_a$，会使得此时圆盘像的半径变为 (1) 中圆盘像的半径的一半？若存在，请给出 $r_a$ 的数值。
    \end{subproblem}
    \begin{subproblem}
        若将圆形光阑移至凸透镜后方 $18\mathrm{cm}$ 处，回答 (2) 中的问题；
    \end{subproblem}
    \begin{subproblem}
        圆形光阑放置在哪些位置时，圆盘像的大小将与圆形光阑的半径有关？
    \end{subproblem}
    \begin{subproblem}
        若将图中的圆形光阑移至凸透镜前方 $6\mathrm{cm}$ 处，回答 (2) 中的问题。
    \end{subproblem}
\end{problem}